\section{Certified finite affine atlases}

\subsection{Strict principal atlases}

The simplest geometric gluing object has one principal overlap for each ordered pair of charts.

\begin{definition}[Strict certified affine atlas]\label{def:strict-atlas}
A strict certified affine atlas consists of:
\begin{enumerate}[label=(\alph*)]
\item a positive integer $r$ and finitely presented rings $A_i$ for $i<r$;
\item elements $f_{ij}\in A_i$ for every $i,j<r$, with $f_{ii}=1$;
\item certified isomorphisms
\[
           \phi_{ij}:(A_i)_{f_{ij}}\xrightarrow{\sim}(A_j)_{f_{ji}};
\]
\item identity witnesses for $\phi_{ii}$ and inverse witnesses for
      $\phi_{ji}\phi_{ij}=\mathrm{id}$;
\item on every triple $(i,j,k)$, a finite common localization and a witness that
\[
                         \phi_{ik}=\phi_{jk}\phi_{ij}
\]
there.
\end{enumerate}
Every ring, element, map, localization, and equality in this definition is a natural-number code
with the certificates described in the preceding sections.
\end{definition}

Externally, the usual gluing theorem for schemes turns this data into a scheme covered by the affine
schemes $\Spec(A_i)$.  Internally, no topological quotient is needed.  A theorem about the glued
object is translated into a finite theorem about the charts and the certified overlap maps.

The strict form already covers projective space: on the chart $X_i\neq0$, the overlap with the chart
$X_j\neq0$ is the principal open where $X_j/X_i$ is invertible.

\subsection{Finite principal-overlap atlases}

Pairwise intersections of affine opens need not always be one principal open.  The practical
extension replaces $f_{ij}$ by a finite family $f_{ij\alpha}$, each giving a principal overlap piece.
The data include:

\begin{itemize}
\item a unit-ideal or relative-cover certificate that the pieces cover the intended intersection;
\item a localization isomorphism on every piece;
\item compatibility witnesses on double refinements;
\item cocycle witnesses on triple refinements.
\end{itemize}

The principal-cover gluing theorem proves that chartwise constructions independent of the chosen
piece descend to the whole intersection.  Refinement coherence follows from uniqueness.  Thus the
additional indexing is finite bookkeeping rather than a new logical principle.

\subsection{The internal geometric language}

A finite atlas supports the following notions without constructing an internal prime spectrum.

\begin{description}[style=nextline]
\item[Regular function.]
A tuple $a_i\in A_i$ with localization equality witnesses
$\phi_{ij}(a_i)=a_j$ on overlaps.  The principal-cover theorem glues such tuples whenever one chart
is itself presented by a finite principal cover.

\item[Morphism.]
A chartwise finite family of presented ring maps in the contravariant direction, together with a
finite target-chart cover and overlap commutativity witnesses.  Relational finite-use composition
may choose among several valid target charts; the induced map on an invariant observation is
choice-free.

\item[Open subobject.]
A finite family of principal elements on each chart with compatibility certificates under overlap
maps.  Union and intersection are computed by concatenation and products of principal elements.

\item[Closed subobject.]
A finite ideal presentation on each chart, stable under localization and equipped with overlap
identities.  A global finitely presented closed subscheme is represented by finitely many local
relation lists.

\item[Point over a presented ring $R$.]
A compatible family of certified ring maps $A_i\to R$ on a selected chart cover, or more commonly a
single chart map with a certificate that its denominators are invertible.  Projective homogeneous
coordinates are a compact specialized presentation of the same data.
\end{description}

These are theorem-relative representations.  A statement requiring only chartwise equations never
needs to enumerate the topological points of the scheme.

\subsection{Atlas morphisms and refinement}

A refinement replaces a chart $A_i$ by a finite principal cover.  Its certificate contains the
unit-ideal relation and the localization maps.  A common refinement of two finite atlases is formed
by taking finite pairwise intersections and then principal refinements.  All index bounds are
primitive recursive in the atlas sizes.

\begin{proposition}[Refinement invariance]\label{prop:refinement-invariance}
Suppose a finite atlas observation is defined by compatible local data and has uniqueness under the
principal-cover equalizer.  Then its value is invariant, with an explicit equality witness, under
finite principal refinement.
\end{proposition}

\begin{proof}
Pull the data to the common refinement.  Both candidate values have equal restrictions on every
refined principal piece.  Apply principal-cover uniqueness chart by chart and then on the outer
finite atlas.  The resulting equality witness is obtained by nested applications of the power-cover
constructor.
\end{proof}

This proposition supplies the practical replacement for proving equality of two constructions on
all open sets.

\subsection{Base change and fiber products}

For presented affine rings, tensor products have explicit presentations.  If
$A\leftarrow R\to B$ are finitely presented, then $A\otimes_R B$ is presented by the combined
variables and relations, together with equations identifying the two images of the generators of
$R$.  The construction is primitive recursive on the finite presentations.

Base-changing an atlas along a presented morphism applies this tensor construction chartwise and
base-changes each localization certificate.  Fiber products of finite atlases are obtained by the
finite product of chart index sets.  Overlap cocycles are inherited by functoriality of localization
and tensor-product maps.  Thus many finite-limit constructions in algebraic geometry fit the same
certificate language.

\subsection{Properties as finite certificates}

Some geometric properties have direct finite witnesses on a chosen atlas:

\begin{itemize}
\item \emph{finite presentation}: already part of every chart and morphism code;
\item \emph{separatedness}: a finite presentation/certificate for each affine diagonal or,
      in common cases, affine pairwise intersections with multiplication-map surjectivity witnesses;
\item \emph{smoothness of a fixed relative dimension}: a finite principal cover on which specified
      Jacobian minors are invertible, together with identities certifying the relations and rank;
\item \emph{finite locally free module of rank $d$}: local bases and inverse matrices on a finite
      principal cover;
\item \emph{proper projective presentation}: a closed homogeneous presentation inside a fixed
      projective space.
\end{itemize}

These certificates are sufficient conditions with soundness proofs.  The framework does not assert
that a uniform primitive-recursive algorithm discovers such certificates for every input
presentation.

\subsection{Conventional semantics}

Externally, Definition~\ref{def:strict-atlas} is interpreted by the standard gluing theorem for
schemes.  The finite ring maps determine open immersions of affine schemes, and the inverse and
cocycle witnesses provide the usual gluing datum.  Every internal chartwise theorem is sound for the
resulting conventional scheme because its algebraic restriction maps are exactly the localization
maps used by the structure sheaf.

A full formal equivalence between the finite-atlas category and an appropriate category of
quasi-compact schemes with chosen atlases would be valuable, but it is not required for individual
iPRA theorems.  Each theorem can state its atlas hypotheses explicitly and prove the desired finite
observation directly.
